\chapter{INTRODUÇÃO}

A digitalização da saúde tornou-se essencial para a otimização do processo, melhoria da segurança da informação e facilitação do acesso a dados médicos. O armazenamento e gerenciamento de documentos médicos são um desafio para clínicas e hospitais, que devem garantir a confidencialidade, integridade e disponibilidade dos dados de seus pacientes. Dado esse cenário, este trabalho propõe o desenvolvimento de uma plataforma de armazenamento de documentos médicos, integrando tecnologias de criptografia, inteligência artificial e assinatura digital para garantir a segurança e a eficiência no gerenciamento dessas informações.

A plataforma permitirá o registro dos pacientes, a declaração de relações e o apego aos seus perfis e a implementação da inteligência artificial para ajudar a criar relatórios e responder perguntas com base em dados registrados. Além disso, a plataforma contará com mecanismos de assinatura digital através de APIs do ClickSign e os recursos de criptografia para garantir a privacidade dos documentos. Esta solução tem como objetivo oferecer um sistema eficaz e seguro para profissionais de saúde e instituições médicas, melhorando digitalmente a acessibilidade e o gerenciamento de documentos médicos.

Neste sentido, a proposta deste trabalho difere dos anteriores por apresentar uma plataforma direcionada ao armazenamento seguro de documentos psicoterapêuticos, com ênfase em saúde mental. A solução integra criptografia de ponta a ponta, mecanismos de assinatura digital com validade jurídica e um modelo de IA treinado localmente para operar em conformidade com os princípios de confidencialidade exigidos pela legislação. Essa abordagem visa preencher uma lacuna observada na literatura nacional, ao combinar privacidade, automação e suporte clínico de forma especializada para a prática psicológica.


\section{Objetivos}

\subsection{Objetivo Geral}
	
Este trabalho tem como objetivo principal a criação de uma plataforma para armazenar de modo seguro, documentos médicos psicoterapêuticos, e por meio de modelos de inteligência artificial possibilitar a geração automatizada de relatórios como laudos psicológicos, pareceres psicológicos, encaminhamentos  entre outros documentos. 

A ideia é que com esta plataforma os profissionais da área tenham um local seguro para armazenar seus documentos evitando extravios e vazamento de informações sensíveis de seus pacientes, além disso, com os documentos salvos em nuvem os profissionais terão mais acessibilidades e eficiência na gestão de seus documentos, e com o uso da Inteligência artificial os profissionais terão maior eficiência na geração de documentos e maior velocidade na obtenção de informações sobre os pacientes cadastrados. 



\subsection{Objetivos Específicos} 

Para um melhor entendimento do objetivo principal deste trabalho, podemos representar nos quatro objetivos específicos a seguir: 

\begin{itemize}
    \item Desenvolver uma plataforma para fazer o armazenamento seguro de documentos psicoterapêuticos.
    \item Disponibilizar a assinatura digital via API externa, para garantir a validade legal dos documentos armazenados.
    \item Utilização de criptografia para trazer mais segurança às informações salvas na plataforma, e evitar o vazamento de informações sensíveis de pacientes.
    \item Implementar um modelo de IA para a geração automatizada de laudos entre outros relatórios terapêuticos, além de extração de informações e resumos de seções.
\end{itemize}
% 
%--------- FIM INTRODUÇÃO------------
%